\section{Introduction}
\label{V:introduction}

An effective divisor can have normalized order tending to zero while
vanishing at the same center in every degree. For example, order one
becomes negligible after division by an unbounded sequence of integers.
Producing a section which is actually nonzero at that center is a stronger
problem. We prove an exact attainment theorem for divisors approaching
the anticanonical class of a smooth projective variety with a smoothly
semipositive anticanonical bundle. We then use generic orbit arcs to
translate this control of orders into statements about characters and
invariant sections.

The broader nonvanishing problem asks whether a nef anticanonical
bundle on a smooth projective variety has a nonzero section in some
positive power.  Numerical effectivity and effectivity in rational
linear equivalence are distinct here.  Lazi\'c, Matsumura, Peternell,
Tsakanikas and Xie prove numerical effectivity for smooth projective
threefolds with nef anticanonical class and develop nonvanishing methods based on twisted
differential forms \cite[Theorem~A and Section~5]{LMPTX2023}.
M\"uller proves anti-log-canonical nonvanishing for projective klt
threefold pairs with nef anti-log-canonical divisor \cite[Corollary~B]{Muller2025}.
Our positivity hypothesis is smooth semipositivity, and our first
result retains both rational linear equivalence and the order at a
prescribed valuation.

Let $X$ be a smooth connected projective complex variety, of dimension
$n>0$, and put $L=-K_X$. Fix a smooth Hermitian metric on $L$ with
semipositive curvature form $\alpha$, normalized to represent $c_1(L)$.
Let $\nu$ be the maximum complex rank of $\alpha$, and let $G_0$ be an
ample line bundle. Define
\begin{equation}\label{V:null-functional}
 \ell(D)=
 \begin{cases}
 D\cdot L^\nu\cdot G_0^{n-\nu-1},&\nu<n,\\
 0,&\nu=n.
 \end{cases}
\end{equation}
This functional is nonnegative on effective divisors. Moreover
$\ell(L)=0$, since $\alpha^{\nu+1}=0$ when $\nu<n$.
A discrete valuation $v$ of $\C(X)$ is \emph{centered on $X$} if its
valuation ring dominates the local ring of a point of $X$. We allow the
zero valuation as well. The value $v(D)$ of a Cartier divisor is the
valuation of a local defining rational function at the center. It is
well defined because units have value zero, and extends rationally to
rational divisors.

\begin{samepage}
\begin{theorem}[Exact attainment of zero order]\label{V:exact}
With $X,L,\alpha,\nu,G_0,\ell$ as above, let $v$ be either a discrete
valuation of $\C(X)$ centered on $X$ or the zero valuation. Let $M$
be any fixed line bundle. Suppose that
effective integral divisors $N_j$ and positive integers $m_j\to\infty$
satisfy
\[
 N_j\sim m_jL+M,\qquad \ell(N_j)=0,\qquad
 \frac{v(N_j)}{m_j}\longrightarrow0.
\]
There is an effective rational divisor $N\sim_{\Q}L$ such that
$v(N)=0$.
\end{theorem}
\end{samepage}

The null intersection condition is the geometric input; the auxiliary
line $M$ is arbitrary. No irregularity or Euler-characteristic hypothesis
is used. The theorem preserves rational linear equivalence, which is
necessary for constructing sections. A limit in the numerical effective
cone would not give its conclusion.

Here is a useful special case. A \emph{generic formal arc} on $X$ is a
morphism $\gamma:\Spec F[[t]]\to X$, for an extension field $F/\C$,
whose fraction point maps to the generic point of $X$. If $-M$ is
pseudoeffective and $N_j\sim m_jL+M$ is effective, nef intersection
gives $0\le\ell(N_j)=\ell(M)\le0$. Thus, if the orders of $N_j$
along $\gamma$ are $o(m_j)$, Theorem~\ref{V:exact} supplies a section
of some positive multiple of $L$ whose order along $\gamma$ is exactly
zero. Corollary~\ref{cor:arc} verifies the centered-valuation and
integrality details, including a point variety and the zero valuation.

\subsection{Why finite generation attains the limit}
The proof first places the principal corrections to the given divisors
in one function field, and then uses finite generation to attain the
limiting class.  Let $k\subset\C(X)$ be the field of rational functions
whose polar divisors have $\ell$-value zero.  It is relatively
algebraically closed.  Two fixed input divisors give representatives
of $L+M/m_j$ on a common finite support, and the principal corrections
between those representatives and $N_j/m_j$ belong to $k$.

The geometric task is to find a model $Y$ of $k$ on which these
corrections lie in a finitely generated divisor section ring.  On a smooth resolution $\pi:Z\to X$ of the rational map to $Y$,
the null-polar field makes the section space of the effective divisor
$K_{Z/X}$ one-dimensional on the generic fiber over $Y$.  The null
intersection also forces the curvature kernel to be vertical on a
nonempty open set.  Section~\ref{V:null-field} verifies these two
inputs to the volume-metric theorem of \emph{Metric descent and rank-preserving
contractions} \cite[Theorem \textup{(Toroidal volume metric)}]{CompanionG}.
It gives an effective rational correction $A_0$ on $Y$ whose pullback
is exceptional over $X$.  The semipositive metric it supplies on $-K_Y+A_0$ has
locally integrable squared frame norms and is smooth and strictly
positive on an ordinary open set.  The correction-ring theorem of
\emph{Integrable metrics and effectivity with controlled boundary}
\cite[Theorem \textup{(Finite generation near a correction divisor)}]{CompanionE}
then gives finite generation in the two rational divisor directions
selected by the input sequence.

The remaining argument uses the actual rational functions representing
sections.  In a finitely generated bigraded ring, a monomial whose
normalized degree approaches a boundary ray uses off-ray generators
in negligible degree.  For each input section, the valuation inequality
selects a monomial with no larger valuation.  One boundary-ray generator
therefore has sufficiently small normalized valuation.  Its divisor
is effective up to the correction $A_0$.  Birational pushforward removes
that correction; centeredness then turns the resulting upper bound on
the valuation into exact equality with zero.  This order of operations
allows the exceptional correction itself to have positive valuation.

The endpoint argument belongs to the study of asymptotic vanishing and
multigraded section rings.  Ein, Lazarsfeld, Musta\c{t}\u{a}, Nakamaye
and Popa explain how finite generation controls asymptotic orders and
base loci \cite[Section~4]{ELMNP2006}, with the corrected polyhedral
argument \cite[Proposition~1.1]{ELMNP2023Erratum}; Corti and Lazi\'c
give a multigraded formulation \cite[Theorem~3.5]{CortiLazic2013}.
Lemma~\ref{lem:endpoint} gives the direct one-ray, one-valuation
argument used here.  The geometric finite-generation input rests on
the adjoint-ring theory of Birkar, Cascini, Hacon and McKernan
\cite{BCHM2010} and Cascini--Lazi\'c \cite{CasciniLazic2012}, in the
form \cite[Theorem~3.2(1)]{CortiLazic2013}.

\subsection{From orders to invariant sections}
For a torus action, the differential of the action gives $-K_X$ its
natural linearization.  A nonzero section can transform by a nontrivial
character, so scalar nonvanishing does not by itself give invariance.
Along the generic orbit of a one-parameter subgroup, an eigensection
of weight $w$ has order $\mu-w$, where $\mu$ is the fiber weight at
the limit.  Exact order zero therefore attains that fiber weight.

M\"uller's equivariant nonvanishing theorem gives invariant
anti-log-canonical sections for projective sub-log-canonical pairs
with semiample anti-log-canonical divisor, under commutative linear
algebraic group actions \cite[Theorem~C]{Muller2025}.  For smooth
varieties with a torus action, we obtain the following result under
smooth semipositivity and a cohomological hypothesis.

\begin{theorem}\label{V:invariant-main}
Let $X$ be smooth connected projective over $\C$, with $-K_X$ smoothly
semipositive and $H^i(X,\OO_X)=0$ for every $i>0$. For every algebraic
torus $T$ acting on $X$, some positive anticanonical multiple has a
nonzero $T$-invariant section for its natural linearization.
\end{theorem}

The proof starts with the source component, which attracts the generic
orbit as the one-parameter subgroup tends to zero.  The attracting
decomposition of Bia\l ynicki-Birula identifies this component
\cite{BB1973}.  The holomorphic fixed-point formula of Atiyah--Bott,
Atiyah--Segal and Atiyah--Singer
\cite{AtiyahBott1968,AtiyahSegal1968,AtiyahSinger1968} supplies
cohomology classes of the corresponding extremal weight.  Semipositive
hard Lefschetz \cite[Theorem~0.1]{DPS2001} turns them into twisted
differential forms.  Saturated determinants, following the method of
Lazi\'c--Peternell \cite[Lemma~4.1]{LP2018} and its anticanonical
adaptation \cite[Lemma~5.1]{LMPTX2023}, give effective
divisors with bounded order along the generic orbit.  Their character
calculation retains the determinant-rank shift.  Theorem~\ref{V:exact}
then supplies actual normalized section weights whose convex hull
contains zero.  A rational convex combination becomes
an invariant product after clearing denominators.  The comparison of
fixed-fiber and section-weight polytopes in the semiample setting is
developed in \cite[Proposition~3.13 and Corollary~3.14]{Muller2025};
here exact attainment supplies the supporting section weights.

Together with the structural inputs described below, the invariant
sections yield the following scalar consequence: for every smooth
connected projective complex variety $X$ with smoothly semipositive
$-K_X$, there is an integer $m>0$ such that
\[
 H^0(X,-mK_X)\ne0.
\]
Sections~\ref{V:isometric-descent} and \ref{V:sec-mrc-descent} give
two descents.  The first uses the compact-monodromy structure theorem
of \emph{Anticanonical nonvanishing from smooth semipositivity}
\cite[Theorem \textup{(Finite cover with compact torus monodromy)}]{CompanionH},
based on the universal-cover structure of
Demailly--Peternell--Schneider and its rational-connectedness refinement
by Campana--Demailly--Peternell \cite{DPS1996,CDP2015}.
The second uses a locally constant maximal rationally connected
fibration, supplied by Matsumura--Wang \cite{MatsumuraWangV3},
following the earlier smooth structure theorem of Cao--H\"oring
\cite{CaoHoring}.  It follows M\"uller's invariant-section reduction
\cite[Theorem~2.2]{Muller2025} and includes the unipotent part of the
commutative algebraic monodromy.  The compact-monodromy argument
uses only the structural and finite-cover norm results of the
nonvanishing companion; its nonvanishing theorem is not an input.

\subsection{Independent methods and their additional outputs}
The exact-order theorem and its invariant application are followed by
two independent constructions.  Their purpose is to retain different
information during the passage from twisted sections to anticanonical
sections.

The first uses nef reduction and extension.  Given bounded orders in
$|m_jL+M|$, with $q(X)=0$ and $-M$ pseudoeffective, it separates a
fixed divisor of zero arc order from a line bundle on a base.  A
fiber-maximum metric makes that line nef; the nef reduction of Bauer
and collaborators \cite{BCEKPRSW2002} removes its numerical null
directions.  Extension of first jets by the theorem of
Cao--Demailly--Matsumura \cite{CDM2017} proves that the remaining base
line is big and yields arbitrarily small normalized arc orders.
Theorem~\ref{V:approximate-arc} records this approximate transfer,
which the exact theorem already strengthens under the same hypotheses.
The independent proof also gives the bigness and jet construction.
Its character application, Proposition~\ref{V:extremal-invariants},
uses $q(X)=0$ and $\chi(X,\OO_X)\ne0$ instead of vanishing of all
higher structure-sheaf cohomology, together with the stated equivariant
very ample linearization.

There is a second ending to this nef-reduction construction.  Starting
with invariant sections of $jL-G$ for unbounded $j$, where $G$ is a
fixed pseudoeffective linearized line bundle, it forms a base from
invariant rational functions with pseudoeffectively controlled poles.
It then extends one nonzero invariant class in a multiplier-ideal
quotient.  Theorem~\ref{V:invariant-conversion} thus removes the fixed
error while keeping invariance throughout.  This ending needs
$q(X)=0$ but does not require the first-jet or weight-polytope steps.
The common tools occupy Section~\ref{V:sec-nef}, the arc and character
applications Sections~\ref{V:sec-transfer} and \ref{V:sec-extremal},
and the invariant conversion Section~\ref{V:sec-invariant-field}.

The other finite-generation construction assumes rational
connectedness and controls an entire collection of divisor classes.
A prime divisor $D$ is controlled when $aL-D$ is pseudoeffective for
some positive integer $a$.  Proposition~\ref{V:controlled-cone} places
the naturally linearized classes of all invariant controlled primes
in one rational polyhedral cone generated by invariant effective
rational divisors.  Its generators need not themselves be controlled.
Here the effectivity companion's weighted-approximation theorem
produces an adjoint boundary on a base; its correction disappears
under birational pushforward.  Finite generation of the full
multigraded ring and then of its invariant subring gives the cone.
Character separation supplies another proof of invariant nonvanishing
in the rationally connected case, as well as the fiberwise input used
in the locally constant MRC descent.  Section~\ref{V:controlled-cone-section}
keeps this weighted metric construction distinct from the unweighted
integrability used for exact attainment.

The proof of exact attainment is organized as follows.
Section~\ref{sec:generator} proves the finite-generator lemma and its
birational application.  Sections~\ref{V:null-field} and
\ref{V:exact-proof} construct the null-polar base and prove
Theorem~\ref{V:exact}; Section~\ref{sec:arc} gives the formal-arc
consequence.  Section~\ref{V:sec-order-weight} establishes the
order/weight identity, and Section~\ref{subsec:valuation-extremal-weights}
proves Theorem~\ref{V:invariant-main}.  The independent methods then
supply the additional conclusions just described, before the paired
monodromy descents finish the scalar nonvanishing consequence.
